\section{Conclusions}
\label{sec:conclusions}

This report applied Bayesian linear regression to model daily bike-sharing demand
as a function of normalised temperature and humidity, using 360 chronologically
subsampled observations from the UCI Bike Sharing Dataset
(\texttt{random\_state = 42}).

\textbf{Likelihood and model specification.} The Normal likelihood was selected
over Poisson on the basis of severe overdispersion (Var/Mean $\approx 853$) and
the approximate symmetry of the empirical rental distribution. The two-predictor
model $\text{cnt}_i \sim \mathcal{N}(\beta_0 + \beta_1\,\text{temp}_i +
\beta_2\,\text{hum}_i,\,\sigma^2)$ was fitted with weakly informative Normal
priors ($\sigma = 10{,}000$) on regression coefficients and a Half-Normal prior
(scale $= 3{,}000$) on the residual standard deviation, ensuring that the
posterior was data-dominated.

\textbf{Posterior inference and frequentist comparison.} NUTS sampling with four
chains produced 8{,}000 posterior draws with $\hat{R} = 1.000$ and bulk
ESS $> 3{,}500$ for all parameters, confirming reliable convergence. Temperature
is the dominant predictor (posterior mean $\hat{\beta}_1 \approx 6{,}784$;
95\% CrI wholly positive), while humidity exerts a credibly negative effect
(posterior mean $\hat{\beta}_2 \approx -2{,}350$; 95\% CrI wholly negative).
Posterior means agree with OLS point estimates to within 1\%, and credible
intervals are nearly indistinguishable in width from the corresponding
frequentist 95\% confidence intervals. This near-coincidence reflects the
dominance of the likelihood under the weakly informative prior at $n = 360$,
while the two frameworks remain conceptually distinct: the Bayesian credible
interval carries a direct probability statement about the parameter, whereas
the frequentist confidence interval is a repeated-sampling construct.

\textbf{Sampler behaviour.} Trace plots showed stable mixing with no trends
or drifts, and all four chains overlapped throughout sampling. Burn-in of
1{,}000 warm-up steps was sufficient for the sampler to reach the stationary
distribution, as confirmed by $\hat{R} = 1.000$ and ESS $> 3{,}500$ for all
parameters.

\textbf{Prior sensitivity.} The sensitivity analysis confirmed that the posterior
is robust to substantial changes in prior specification. Together with the close
agreement between the Bayesian and frequentist estimates, this indicates that
360 observations provide sufficient information to overwhelm moderate variations
in prior beliefs.

\textbf{Sequential Bayesian updating.} The extension task partitioned the data
into four chronological blocks of 90 observations and cascaded posteriors as
priors via moment matching. Several findings emerged. First, the temperature
coefficient $\beta_1$ converged reliably across blocks: its posterior standard
deviation fell by 40\% from Block 1 (403.7) to Block 4 (241.6), and its
credible interval remained wholly positive throughout, confirming the robustness
of the temperature signal. Second, the humidity coefficient $\beta_2$ proved
difficult to identify within individual 90-observation windows: credible intervals
spanned zero in three of four blocks, and the posterior mean changed sign between
blocks. This instability was attributed to two factors --- small-sample noise
insufficient to resolve a weaker signal, and genuine seasonal variation in the
local temperature--humidity--rental relationship --- rather than
multicollinearity, which was ruled out by between-predictor correlations no
greater than 0.246 in any block. Third, the residual standard deviation $\sigma$
converged from an initial estimate of 652 in Block 1 to 1{,}497 in Block 4,
closely matching the full-data estimate and demonstrating that the likelihood
eventually overrides an initially displaced prior. These results highlight an
important asymmetry in sequential updating: dominant signals are reliably learned
block by block, while weaker signals require larger individual segments or more
informative starting priors to avoid entrenchment of early estimation errors.

\textbf{Limitations and future directions.} The two-predictor model deliberately
omits season, weekday, and holiday effects that account for the remaining 60\%
of unexplained variance ($R^2 = 0.405$). Future work could incorporate these
covariates, adopt a Negative Binomial likelihood to address residual
overdispersion, or employ a hierarchical structure to share information across
seasonal strata. For the sequential analysis specifically, larger block sizes
or more informative priors derived from domain knowledge about humidity effects
would mitigate the small-sample instability observed for $\beta_2$.
